
\documentclass{article}
\usepackage{colortbl}
\usepackage{makecell}
\usepackage{multirow}
\usepackage{supertabular}

\begin{document}

\newcounter{utterance}

\centering \large Interaction Transcript for game `codenames', experiment `unambiguous', episode 4 with reflexion{-}llama{-}3.1{-}8b{-}instruct.
\vspace{24pt}

{ \footnotesize  \setcounter{utterance}{1}
\setlength{\tabcolsep}{0pt}
\begin{supertabular}{c@{$\;$}|p{.15\linewidth}@{}p{.15\linewidth}p{.15\linewidth}p{.15\linewidth}p{.15\linewidth}p{.15\linewidth}}
    \# & $\;$A & \multicolumn{4}{c}{Game Master} & $\;\:$B\\
    \hline

    \theutterance \stepcounter{utterance}  
    & & & \multicolumn{2}{p{0.3\linewidth}}{
        \cellcolor[rgb]{0.9,0.9,0.9}{
            \makecell[{{p{\linewidth}}}]{
                \texttt{\tiny{[GM$|$GM]}}
                {'hidden': {'team': ['stairway'], 'innocent': ['stewardess', 'blimp', 'appetizer', 'pharmacist', 'prom', 'geisha', 'shaman'], 'opponent': [], 'assassin': ['cowardice']}, 'revealed': {'team': {'team': [], 'innocent': [], 'opponent': [], 'assassin': []}, 'opponent': {'team': [], 'innocent': [], 'opponent': [], 'assassin': []}}}
            }
        }
    }
    & & \\ \\

    \theutterance \stepcounter{utterance}  
    & & \multicolumn{4}{p{0.6\linewidth}}{
        \cellcolor[rgb]{0.9,0.9,0.9}{
            \makecell[{{p{\linewidth}}}]{
                \texttt{\tiny{[P1$\langle$GM]}}
                \texttt{You have attempted this same task before and did not succeed. These are the plans you wrote after those attempts.} \\
\texttt{Trial \#1: New plan: } \\
\\ 
\texttt{In my previous attempts, I focused on providing clues that were semantically related to multiple words, but this led to my teammate guessing opponent words, which ended their turn. I should have been more selective in choosing the words to relate to the clue. Next time, I will provide a clue that is more closely related to a single word or a smaller group of words, increasing the chances of my teammate guessing a team word correctly. I will also pay closer attention to the words that have been removed from the lists and avoid providing clues that might lead to guessing opponent words.} \\
\\ 
\texttt{Trial \#2: New plan: In my previous attempts, I focused on providing clues that were semantically related to multiple words, but this led to my teammate guessing opponent words, which ended their turn. I should have been more selective in choosing the words to relate to the clue. Next time, I will provide a clue that is more closely related to a single word or a smaller group of words, increasing the chances of my teammate guessing a team word correctly. I will also pay closer attention to the words that have been removed from the lists and avoid providing clues that might lead to guessing opponent words. Specifically, I will start by examining the team words and looking for a common thread or a unique characteristic that can be used as a clue.} \\
\\ 
\texttt{Trial \#3: New plan: In my previous attempts, I focused on providing clues that were semantically related to multiple words, but this led to my teammate guessing opponent words, which ended their turn. I should have been more selective in choosing the words to relate to the clue. Next time, I will provide a clue that is more closely related to a single word or a smaller group of words, increasing the chances of my teammate guessing a team word correctly. I will also pay closer attention to the words that have been removed from the lists and avoid providing clues that might lead to guessing opponent words. Specifically, I will start by examining the team words and looking for a common thread or a unique characteristic that can be used as a clue. I will also consider the distractor words and try to find a clue that is more closely related to the team words than to the distractor words.} \\
\\ 
\\ 
\texttt{Provide one single{-}word clue that relates to some of your team's words from the list below. You can choose to how many and to which words your clue relates to.} \\
\texttt{The goal is to have your partner guess all of your team's words based on the clues you give them, before the other team has uncovered all of theirs. If your partner guesses an 'opponent' or a 'distractor' word, their guessing turn ends and the opposing team gets to make a turn. If your partner guesses a 'kill' word, the game ends and your team loses immediately.} \\
\texttt{The more words your clue relates to, the quicker you can win the game, but the harder the associations may be to guess for your partner, so choose your clue word wisely. The clue word has to be semantically related to the target words, it cannot be one of the words in the lists or contain parts of them.} \\
\\ 
\texttt{Always give your single{-}word clue and your comma{-}separated list of related target words in the following format and make your answers as short as possible, never include any other text than is required in this form:} \\
\\ 
\texttt{CLUE: <WORD>} \\
\texttt{TARGETS: <TARGETS>} \\
\\ 
\texttt{Your team words are: stairway.} \\
\texttt{Your opponent's team words are: .} \\
\texttt{Distractor words are: stewardess, blimp, appetizer, pharmacist, prom, geisha, shaman.} \\
\texttt{The kill words are: cowardice.} \\
            }
        }
    }
    & \\ \\

    \theutterance \stepcounter{utterance}  
    & \multicolumn{4}{p{0.6\linewidth}}{
        \cellcolor[rgb]{0.9,0.9,0.9}{
            \makecell[{{p{\linewidth}}}]{
                \texttt{\tiny{[P1$\rangle$GM]}}
                \texttt{CLUE: ascent} \\
\texttt{TARGETS: stairway} \\
            }
        }
    }
    & & \\ \\

    \theutterance \stepcounter{utterance}  
    & & & \multicolumn{2}{p{0.3\linewidth}}{
        \cellcolor[rgb]{0.9,0.9,0.9}{
            \makecell[{{p{\linewidth}}}]{
                \texttt{\tiny{[GM$|$GM]}}
                \texttt{ascent} \\
            }
        }
    }
    & & \\ \\

    \theutterance \stepcounter{utterance}  
    & & & \multicolumn{2}{p{0.3\linewidth}}{
        \cellcolor[rgb]{0.9,0.9,0.9}{
            \makecell[{{p{\linewidth}}}]{
                \texttt{\tiny{[GM$|$GM]}}
                ['stairway']
            }
        }
    }
    & & \\ \\

    \theutterance \stepcounter{utterance}  
    & & & \multicolumn{2}{p{0.3\linewidth}}{
        \cellcolor[rgb]{0.9,0.9,0.9}{
            \makecell[{{p{\linewidth}}}]{
                \texttt{\tiny{[GM$|$GM]}}
                {'word': 'stairway', 'assignment': 'team'}
            }
        }
    }
    & & \\ \\

    \theutterance \stepcounter{utterance}  
    & & & \multicolumn{2}{p{0.3\linewidth}}{
        \cellcolor[rgb]{0.9,0.9,0.9}{
            \makecell[{{p{\linewidth}}}]{
                \texttt{\tiny{[GM$|$GM]}}
                \texttt{opponent has won} \\
            }
        }
    }
    & & \\ \\

    \theutterance \stepcounter{utterance}  
    & \multicolumn{4}{p{0.6\linewidth}}{
        \cellcolor[rgb]{0.95,0.95,0.95}{
            \makecell[{{p{\linewidth}}}]{
                \texttt{\tiny{[P2$\rightarrow$P1]}}
                \texttt{Post{-}episode Reflexion (not sent to GM):} \\
\texttt{New plan: In my previous attempts, I focused on providing clues that were semantically related to multiple words, but this led to my teammate guessing opponent words, which ended their turn. I should have been more selective in choosing the words to relate to the clue. Next time, I will provide a clue that is more closely related to a single word or a smaller group of words, increasing the chances of my teammate guessing a team word correctly. I will also pay closer attention to the words that have been removed from the lists and avoid providing clues that might lead to guessing opponent words. Specifically, I will start by examining the team words and looking for a common thread or a unique characteristic that can be used as a clue. I will also consider the distractor words and try to find a clue that is more closely related to the team words than to the distractor words. Additionally, I will pay closer attention to the kill words and try to avoid providing clues that might lead to guessing them, as this would end the game immediately.} \\
            }
        }
    }
    & & \\ \\

\end{supertabular}
}

\end{document}
